\section{Actual original Gauss-point norms on a literal test configuration}

\begin{lemma}\label{lem:actual-trivial-polynomial-gauss-bound}
\lean{EnvelopeProperty.Berkovich.Spectrum.gauss_val_le_one, EnvelopeProperty.Berkovich.Spectrum.trivialGaussValuation, EnvelopeProperty.Berkovich.Spectrum.trivialGaussValuation_X, EnvelopeProperty.Berkovich.Spectrum.trivialGaussValuation_le_one}
\leanok
For any original field $k$ and $0<r\leq1$, the actual polynomial Gauss
valuation extending the original trivial valuation is bounded by one on
the entire original $k[t]$ and takes the original parameter to $r\ne0$.
\uses{spectrum:gauss-point}
\end{lemma}
\begin{proof}
The coefficient maximum formula bounds every coefficient value by one
and every nonnegative power of $r$ by one. Each term, and hence the
finite maximum including the zero polynomial, is at most one. The
actual Gauss parameter formula gives its value exactly $r$, which is
nonzero by positivity. No boundedness or parameter-value certificate
is added to the original data.
\end{proof}

\begin{definition}\label{def:native-test-configuration-gauss-power-norms}
\lean{EnvelopeProperty.NativeTestConfiguration.TestConfiguration.gaussPowerBaseIdealNormalizedRoot, EnvelopeProperty.NativeTestConfiguration.TestConfiguration.gaussPowerSectionNorm}
\leanok
Let $C$, its original invertible total line $M$, the original ample
fibre line and the original line-action/fibre-one isomorphisms be as in
the preceding actual test-configuration theorems. At each original native
point $x$ of $X$ extending the trivial valuation and each $0<r\leq1$,
evaluate the actual entire base ideal $I_m$ and each actual original
global section of $M^m$ at the constructed original point $y_{C,x,r}$.
The normalized root is $\|I_m(y_{C,x,r})\|^{1/m}$, using the real root
convention also at zero norms and at degree zero.
\uses{thm:native-test-configuration-gauss-point-base,lem:actual-trivial-polynomial-gauss-bound,lem:native-base-ideal-root-multiples}
\end{definition}
\begin{proof}
The genuine original point has the entire prescribed polynomial Gauss
base by the preceding whole-point theorem. The first lemma gives the
bounded base needed for the actual proper model norms. Source
projectivity supplies the actual properness of the original projection.
Thus these are the actual original proper-model ideal and line norms;
neither a point nor a norm matching certificate is supplied.
\end{proof}

\begin{theorem}\label{thm:native-test-configuration-actual-gauss-norm-supremum}
\lean{EnvelopeProperty.NativeTestConfiguration.TestConfiguration.exists_gauss_normalized_supremum}
\leanok
At every original $x$, these actual normalized roots converge to a
positive real $b(x,r)\leq1$, which is their least upper bound over
\emph{all positive original degrees}. Early zero ideals are included.
\uses{def:native-test-configuration-gauss-power-norms,thm:native-test-configuration-normalized-supremum}
\end{theorem}
\begin{proof}
The boundedness and nonzero parameter value are derived from the first
lemma, and the relative original Gauss point is constructed from $x$.
Apply the original all-positive-degree supremum theorem at that actual
point. Its positive root limit and exact least-upper-bound assertion
are precisely those for the original norms just defined. Degree zero
is excluded from the supremum, exactly as in that theorem.
\end{proof}

\begin{theorem}\label{thm:native-test-configuration-actual-gauss-finite-norm-family}
\lean{EnvelopeProperty.NativeTestConfiguration.TestConfiguration.exists_eventual_finite_gauss_power_section_norms}
\leanok
One degree threshold works for every original $x$. For each degree
beyond it, one finite family of original total-space global sections,
chosen before $x$, computes the actual entire base-ideal norm at every
$y_{C,x,r}$ as the maximum of their actual model line norms. This
maximum is positive and at most one.
\uses{def:native-test-configuration-gauss-power-norms,thm:native-test-configuration-eventual-finite-positive-norm-family}
\end{theorem}
\begin{proof}
The actual polynomial Gauss valuation is fixed before $x$. Apply the
original finite-family theorem with its derived boundedness and nonzero
parameter value. It supplies one threshold and, at each later degree,
one actual global-section family valid at every original point over
this valuation. Specialize its whole assertion to the actually
constructed $y_{C,x,r}$. This preserves the family before the point and
gives the exact original ideal-norm identity and both bounds.
Identification of these actual norms with the weighted Laurent-section
FS metric remains a subsequent, separate geometric comparison.
\end{proof}
